% ---------------------------------------------------------------------------
% Metodología
% ---------------------------------------------------------------------------
% Mínimo dos páginas por capítulo (norma 2.6).
%
% Este capítulo debe DECLARAR las decisiones de diseño, no solo describir el
% sistema: son las que hacen reproducible el experimento del capítulo 4. Las que
% ya están tomadas se listan abajo con su fecha; el seguimiento completo, con las
% que siguen abiertas, está en la sección 8 de study_materials/guia_del_proyecto.md
% (documento de trabajo, no versionado).
%
% DECIDIDO 24/08/2026 — hay que declararlo en el capítulo:
% DECIDIDO   - Idioma de las peticiones: ESPAÑOL, contra un corpus en inglés.
% DECIDIDO     Implica recuperación CRUZADA entre idiomas, y eso condiciona la
% DECIDIDO     elección del modelo de embeddings. Declararlo explícitamente: es
% DECIDIDO     una restricción del diseño, no un supuesto tácito.
% DECIDIDO   - Extracción del documento: pdftotext con la opción -layout, que
% DECIDIDO     preserva la alineación de columnas de las tablas de registros.
% DECIDIDO     Verificado sobre las 39 secciones del datasheet (707 de sus 767 páginas).
% DECIDIDO     Limitación conocida y acotada: las figuras y el sombreado por
% DECIDIDO     dispositivo no sobreviven.
% DECIDIDO   - Formato de documentación de los periféricos: plantilla de seis
% DECIDIDO     secciones en project/perifericos/PLANTILLA.md. Es a la vez corpus
% DECIDIDO     del sistema y patrón de referencia para medir fidelidad.
%
% DECIDIDO 26/08/2026 — divergencias entre el plan oficial y lo que se hace.
% DECIDIDO   Hay que declararlas y justificarlas; no basta con describir el sistema.
% DECIDIDO   - ENTORNO DE DESARROLLO. El objetivo 2 del plan CCT-002-2026 compromete
% DECIDIDO     "MPLAB X IDE, compilador XC8 y MPLAB Code Configurator". De los tres,
% DECIDIDO     DOS se cumplen tal cual: XC8 v3.10 compila los ejemplos, y MCC genera
% DECIDIDO     la configuración de los proyectos (config.mcc/). MPLAB X IDE NO está
% DECIDIDO     instalado: se usa Visual Studio Code con las 13 extensiones de
% DECIDIDO     Microchip, entre ellas mplab-code-configurator.
% DECIDIDO     Argumento para el capítulo: son dos frontales del MISMO toolchain de
% DECIDIDO     Microchip. El compilador y el generador de configuración son idénticos;
% DECIDIDO     lo que cambia es el editor y el formato de proyecto —.vscode/*.mplab.json
% DECIDIDO     más config.mcc/, en vez de un .X—. El objetivo 2 no nombra la placa
% DECIDIDO     ("la tarjeta de evaluación correspondiente"), así que la Curiosity Nano
% DECIDIDO     no genera ninguna divergencia.
% DECIDIDO     El plan NO se corrige: la diferencia se explica aquí.
% DECIDIDO   - MCP. El título del plan compromete "RAG y MCP", pero ninguno de sus
% DECIDIDO     siete objetivos menciona MCP. Se resolvió tratándolo como objeto de
% DECIDIDO     estudio: el capítulo 2 tiene una sección propia sobre el protocolo.
% DECIDIDO     Este capítulo describirá el despliegue del recuperador como servidor
% DECIDIDO     MCP sabiendo que los objetivos del plan no lo listan.
%
% DECIDIDO 02/09/2026 — dos modalidades de MCC, cada una en su papel.
% DECIDIDO   Los seis .mc3 de los ejemplos de Microchip son MCC Classic (core 4.65)
% DECIDIDO   y sirven de corpus de contraste. Los proyectos propios (P01–P03 y los
% DECIDIDO   que sigan) se generan con MCC Melody (core 5.9.1): la extensión de
% DECIDIDO   MPLAB para VS Code solo ofrece Melody, así que el patrón de referencia
% DECIDIDO   nace de proyectos Melody. La fidelidad se contrasta sobre los valores
% DECIDIDO   escritos en los registros, que no dependen de la modalidad. El
% DECIDIDO   capítulo 2 (§2.3) ya lo declara; aquí hay que decir cómo se compara.
% DECIDIDO   - Documentación de periféricos: los .md de project/perifericos/ se
% DECIDIDO     escriben EN LOTE al preparar la ingesta (Fase 4), no al cerrar cada
% DECIDIDO     proyecto. El entregable de la Fase 3 es el proyecto validado en placa
% DECIDIDO     con su FICHA.md; gpio.md (25/08) fue la validación de la plantilla.
%
% PENDIENTE de decidir antes de cerrar este capítulo:
% PENDIENTE   - Modelo de embeddings. Se elegirá MIDIENDO, no por criterio, y
% PENDIENTE     el protocolo está escrito de antemano en
% PENDIENTE     project/banco/PROTOCOLO.md. El capítulo debe declarar: qué
% PENDIENTE     cuenta como acierto (la página de origen), las dos métricas
% PENDIENTE     (recall@k y MRR), el umbral de significación fijado antes de
% PENDIENTE     medir (10 % en recall@5), el criterio de desempate, y que el
% PENDIENTE     corpus medido es solo el datasheet en inglés, que es el caso
% PENDIENTE     MAS DURO para la recuperación cruzada. Reportar la comparación
% PENDIENTE     completa y el análisis de los desacuerdos, no solo el ganador.
% PENDIENTE     Nota de coste: toda cifra de velocidad va con el estado del
% PENDIENTE     procesador declarado; en «powersave» sale 3-5 veces baja.
% PENDIENTE   - Granularidad del fragmento.
% PENDIENTE   - Modelo de lenguaje: API o local, y su temperatura. Con
% PENDIENTE     temperatura no nula hay que declarar el número de repeticiones.
% PENDIENTE   - Límite de iteraciones del lazo de corrección y criterio de
% PENDIENTE     convergencia. Sin ellos, la tasa de éxito no es comparable.
% PENDIENTE   - Tamaño del conjunto de consultas de prueba.

\chapter{Metodología}

% Contenido pendiente.
